% WHAT A STREAM IS OPENED ON.
%
% \openout takes a file name, and a name with NO EXTENSION is opened on
% that name and `.tex':
%
%   \openout3=zzowtest        \openout3 = `zzowtest.tex'.
%   \openout4=zzowtest2.dat   \openout4 = `zzowtest2.dat'.
%   \openout5=zzow.test3.x    \openout5 = `zzow.test3.x'.
%
% The file on disk is named that way too, not just the log line, so that an
% \input or an \openin of the bare name finds it again. An extension is a
% dot in the last part of the name, whatever follows it.
%
% HSTeX opened the bare name. That is invisible until a run writes a file
% and reads it back, which trip does twice: line 94 writes tripos with
% `\openout10=tr\romannumeral1 \gobble\newcs pos', and lines 211 and 424
% read it with \input.
%
% Found and left open: the reference writes ` [the bare name is found
% again]' on a line of its own and `(./zzowtest.tex)' on the next, where
% HSTeX writes both where the last line left off. The reference keeps TWO
% column counters -- one for the terminal and one for the log -- and a
% \message and an \input break their line by the TERMINAL's, which the
% \openout notes moved without moving the log's. HSTeX keeps one.
%
% This probe leaves three files behind; run it where that is welcome.
\catcode`\{=1 \catcode`\}=2
\immediate\openout3=zzowtest
\immediate\write3{one}
\immediate\closeout3
\immediate\openout4=zzowtest2.dat
\immediate\write4{two}
\immediate\closeout4
\immediate\openout5=zzow.test3.x
\immediate\write5{three}
\immediate\closeout5
\message{[the bare name is found again]}
\input zzowtest
\message{[done]}
\end
